\section{Исследование влияния полосы обзора (Span)}

В данном разделе исследуется влияние полосы обзора (\textit{Span}) анализатора спектра на визуальное отображение спектрального пика гармонического сигнала при фиксированной центральной частоте настройки $\text{Center Freq} = 1{,}0\text{ ГГц}$ и неизменной полосе пропускания фильтра промежуточной частоты $\text{RBW} = 300\text{ кГц}$.

\subsection{Результаты эксперимента}
Исходное значение полосы обзора, полученное после автоматической настройки, составляло $\text{Span} = 10\text{ МГц}$. Для анализа влияния масштаба по частоте значение полосы обзора было увеличено в 5 раз (до $50\text{ МГц}$), а затем уменьшено в 5 раз относительно исходного (до $2\text{ МГц}$). Значение полосы фильтра $\text{RBW}$ принудительно фиксировалось на исходном уровне 300 кГц. Полученные спектрограммы приведены на рис.~\ref{fig:span_50mhz} и рис.~\ref{fig:span_2mhz}.

\begin{figure}[H]
    \centering
    \begin{minipage}[b]{0.48\textwidth}
        \centering
        % Замените fig4 на имя вашего файла
        \includegraphics[width=\textwidth]{fig4}
        \caption{Полоса обзора увеличена в 5 раз ($\text{Span} = 50\text{ МГц}$)}
        \label{fig:span_50mhz}
    \end{minipage}
    \hfill
    \begin{minipage}[b]{0.48\textwidth}
        \centering
        % Замените fig5 на имя вашего файла
        \includegraphics[width=\textwidth]{fig5}
        \caption{Полоса обзора уменьшена в 5 раз ($\text{Span} = 2\text{ МГц}$)}
        \label{fig:span_2mhz}
    \end{minipage}
\end{figure}

    
    При изменении параметра $\text{Span}$ изменяется масштаб горизонтальной оси экрана (цена одного деления координатной сетки):

    При увеличении $\text{Span}$ до $50\text{ МГц}$ весь экран отображает в 5 раз более широкий диапазон частот, цена деления сетки увеличивается в 5 раз, из-за чего неизменный по частоте отклик занимает в 5 раз меньше экранных пикселей и визуально кажется узким шпилем;

    При уменьшении $\text{Span}$ до $2\text{ МГц}$ масштаб сетки сильно растягивается (цена деления уменьшается в 5 раз), и тот же самый неизменный частотный отклик занимает значительно большую долю горизонтальной шкалы дисплея, визуально расширяясь.

\textbf{Изменение формы и ширины спектрального пика при изменении диапазона (Span):}

В ходе эксперимента было установлено, что визуальная (экранная) ширина пика изменяется обратно пропорционально ширине полосы обзора (\texttt{Span}).
\begin{itemize}
    \item При увеличении \texttt{Span} в 5 раз (с 10~МГц до 50~МГц) при жестко зафиксированной полосе \texttt{RBW} ($300$~кГц) визуальная ширина пика пропорционально уменьшилась в 5 раз, и сигнал принял форму узкой «иглы».
    \item При уменьшении \texttt{Span} в 5 раз (с 10~МГц до 2~МГц) пик визуально расширился в 5 раз, детально отобразив на экране колоколообразную форму.
\end{itemize}

\textbf{Анализ аппаратных причин и вывод:}

\begin{itemize}
    \item \textbf{Изменение диапазона отображаемых частот} (\texttt{Span}) происходит за счет изменения параметров \textbf{генератора пилообразного напряжения} (или заменяющей его цифровой схемы подсчета импульсов опорного генератора). Этот блок задает пределы перестройки частоты \textbf{гетеродина} и одновременно осуществляет развертку индикатора по горизонтали (ось~X).
    \item \textbf{Влияние полосы фильтра \texttt{RBW} на ширину пика:} Ширина и форма отображаемого спектрального пика для чистого гармонического сигнала зависят от амплитудно-частотной характеристики (АЧХ) и полосы пропускания \textbf{фильтра разрешения (\texttt{RBW})} в тракте последнего усилителя промежуточной частоты (УПЧ)]. Анализатор выводит на экран не бесконечно тонкую дискретную линию, а форму гауссовой кривой этого фильтра. 
\end{itemize}

\textbf{Вывод:} 
Генератор пилообразного напряжения определяет масштаб оси частот и ширину обозреваемого диапазона, а фильтр \texttt{RBW} жестко задает абсолютную частотную ширину отображаемого пика. Изменение параметра \texttt{Span} работает исключительно как визуальное масштабирование (зум), сжимая или растягивая неизменную АЧХ фильтра \texttt{RBW} на экране дисплея.